\documentclass{article}

\usepackage{amssymb}
\usepackage{amsthm}
\usepackage{amsmath}
\usepackage{geometry}
\usepackage{mathrsfs}
\usepackage{stmaryrd}
\usepackage{tikz-cd}

\usepackage[linktoc=all, hidelinks]{hyperref}

\newcommand{\CC}{\mathbb{C}}
\newcommand{\FF}{\mathbb{F}}
\newcommand{\NN}{\mathbb{N}}
\newcommand{\QQ}{\mathbb{Q}}
\newcommand{\RR}{\mathbb{R}}
\newcommand{\ZZ}{\mathbb{Z}}
\newcommand{\id}{\mathrm{id}}
\newcommand{\op}[1]{#1^{\text{op}}}
\newcommand{\powerset}[1]{\mathcal{P}(#1)}

\DeclareMathOperator{\res}{res}
\DeclareMathOperator{\Hom}{Hom}

\newcommand{\Set}{\mathbf{Set}}

\renewcommand\labelenumi{(\alph{enumi})}

\theoremstyle{definition}
\newtheorem{theorem}{Theorem}[section]
\newtheorem{proposition}[theorem]{Proposition}
\newtheorem{definition}[theorem]{Definition}
\newtheorem{lemma}[theorem]{Lemma}
\newtheorem{corollary}[theorem]{Corollary}
\newtheorem{example}{Example}[theorem]

\newtheorem{exercise}{Exercise}[subsection]

\renewcommand{\theexercise}{\thesubsection.\Alph{exercise}}


\title{FOAG The Rising Sea exercises}
\author{bpaul}
\date{\today}

\begin{document}
\maketitle

\section{Category theory}
might do later we'll see

\section{Sheaves}

\subsection{Motivating example: The sheaf of smooth functions}

\begin{exercise}
    Show that \(\mathfrak{m}_p\), the ideal of germs vanishing at \(p\), is the only maximal ideal of \(\mathscr{O}_p\).
    (Hint: show that every element of \(\mathscr{O}_p \setminus \mathfrak{m}_p\) is invertible.)
\end{exercise}

\begin{proof}
    Let \((f, U)\) be represent a germ that is non-zero at \(p\).
    Then there's some open neighbourhood of \(f(v)\) that doesn't contain zero, and the preimage of that neighbourhood is open.
    Call it \(V\).
    Then the restriction \((f |_{U \cap V}, U \cap V)\) is non-zero everywhere and we can define an inverse!
    Hence every element of \(\mathscr{O}_p \setminus \mathfrak{m}_p\) is invertible.
    Any proper ideal must contain no units, and must therefore be contained in \(\mathfrak{m}_p\).
\end{proof}

\begin{exercise}
    Show that \(\mathfrak{m}_p / \mathfrak{m}_p^2\) is the cotangent space at \(p\).
\end{exercise}

\begin{proof}
    skipped (idk DG) something about killing non-linear terms idk
\end{proof}

\subsection{Definition of sheaf and presheaf}
\begin{exercise}
    Show that presheafs are the same thing as contravariant functors from the poset category of open sets to \(\Set\).
\end{exercise}

\begin{proof}
    Associating an \(\mathscr{F}(U)\) to each open \(U\) is defining the map on objects.
    Defining restriction maps \(\res_{V,U} : \mathscr{F}(V) \to \mathscr{F}(U)\) is defining the map on morphisms \(U \hookrightarrow V\).
    Requiring \(\res_{U, U} = \id_{\mathscr{F}(U)}\) is the identity axiom of a functor, and requiring \(\res_{W, V} \circ \res_{V, U} = \res_{W, U}\) is the composition axiom of a functor.
\end{proof}

\begin{exercise}
    Show the following are presheaves on \(\CC\) but not sheaves: (a) bounded functions, (b) holomorphic functions admitting a holomorphic square root.
\end{exercise}
\begin{proof}
    These examples are both presheaves because restriction is function restriction.

    Bounded functions fail gluability since we can take infinite unions.
    For example, take \(U_i = [-1, 1] \times [i, i+1] \subseteq \CC\) with \(f(z) = |z|\).
    This function is bounded on all \(U_i\) but is not bounded an their union.

    Holomorphic functions with holomorphic square roots fail gluability because we can glue over branch cuts.
    For example, we could take the identity function on three thirds of the complex plane with overlap that glue to the whole space, and so each function has a holomorphic square root, but their union does not.
\end{proof}

\begin{exercise}
    What limit expresses the identity and gluability axioms?
\end{exercise}
\begin{proof}
    Gluability and identity can be expressed as, for any open cover \(U_i\) of \(U\), saying that
    \(\begin{tikzcd}[cramped, sep=small] \mathscr{F}(U) \arrow[r] & \prod_i \mathscr{F}(U_i) \arrow[r, shift left=0.75ex] \arrow[r, shift right=0.75ex] & \prod_{i, j} \mathscr{F}(U_i \cap U_j) \end{tikzcd}\)
    is an equaliser diagram.
    We also would want that the cone map is restriction maps, or equivalently that \(\mathscr{F}\) preserves colimits-to-limits (it is contravariant) of this form.
    Then the identity axiom is the fact that \(\mathscr{F}(U)\) is a subobject of the product of \(\mathscr{F}(U_i)\) and gluability is the fact that elements are (can be) represented by tuples that equalise.
\end{proof}

\begin{exercise}
    \begin{enumerate}
        \item
            Verify that the sheaves of smooth functions, continuous functions, real-analytic functions and real-valued functions on a manifold are sheaves.
        \item
            Show that real-valued continuous functions on a topological space \(X\) form a sheaf.
    \end{enumerate}
\end{exercise}
\begin{proof}
    \begin{enumerate}
        \item
            dhgmbwidgmhubwidhgbmuidhgmbw i don't wanna what the heck is a manifold
        \item
            See exercise \ref{ex22f}
    \end{enumerate}
\end{proof}

\begin{exercise}
    Let \(\mathscr{F}(U)\) be the set of maps \(U \to S\) that are locally constant (every point \(p \in U\) has a neighbourhood that the map is constant on).
    Show this forms a sheaf on \(X\).
\end{exercise}
\begin{proof}
    A locally constant map \(f : U \to S\) is precisely a continuous map endowing \(S\) with the discrete topology.
    The neighbourhood of \(p\) will be \(f^{-1}(\{f(p)\})\).
    By exercise \ref{ex22f}, this forms a sheaf!
\end{proof}

\begin{exercise}\label{ex22f}
    Show the sheaf of continuous maps into some space \(Y\) is a sheaf on \(X\).
\end{exercise}
\begin{proof}
    This functor is \(\Hom(\iota(-), Y)\) where \(\iota\) is the inclusion from open sets to topological spaces.
    Since \(\iota\) preserves cover coequalisers, the composite sends cover coequalisers to the correct equalisers.
\end{proof}
\end{document}
